% !PRINTED VII
\weberpage{VII}

\noindent to mention, as far as possible, the most important sources and literary aids at the appropriate places.

It is a pleasant duty for me to express my thanks here to many a friend and colleague who has taken a lively and active interest in the progress of my work. These thanks go first of all to my friend \gesperrt{Dedekind} for his faithful help with the proofreading. And even though he has exercised no influence on the plan and the execution of my work, I should nevertheless not like to leave unmentioned something from many years ago. Through a set of notes of a lecture course on higher algebra, in particular on the theory of \gesperrt{Galois}, which he gave in Göttingen in the winter of 1857/58, I gained an even livelier interest in this theory. Until then it had probably never been presented at our universities with such completeness.

I must also make mention here of the manifold stimulation and instruction that I owe to my friend and colleague F.~\gesperrt{Klein}. He has followed the progress of the work with the liveliest interest, and his expert advice, always most readily given, has been of great influence on many parts of the book.

I cannot name here all the colleagues and friends who, as the reader will find in the relevant places, have supported me with their advice. But I must still make mention here of Messrs.\ E.~\gesperrt{Hess} in Marburg, Fr.~\gesperrt{Meyer} in Clausthal and R.~\gesperrt{Fricke} in Braunschweig. Through their expert and careful execution of the laborious correction of the proof sheets, they have furthered the accuracy and correctness of the text.

Finally, my thanks go to the publishing house, which, by readily accommodating my wishes and by its care in printing and production, has contributed substantially to the success of the whole.

\bigskip

\gesperrt{Göttingen}, November 1894.

\medskip

\begin{flushright}
\textbf{\gesperrt{The Author}.}
\end{flushright}
